% Related work in detail: Reproducible Science beside earlier work by layer, and the rows Section~\ref{sec:related}
% treats briefly or not at all. Every sentence stating what a source does rests on a quotation
% pinned in paper/prior_art/claims/<key>.yaml; `citations verify --claims paper/prior_art/claims`
% resolves each one.
\section{Related Work in Detail}
\label{sec:related-detail}

Table~\ref{tab:related-layers} places Reproducible Science beside earlier work by the layer each system
addresses. Section~\ref{sec:related} treats claim-to-location binding, repository auditing,
artifact-level verification, re-execution and entailment, and the paragraphs below take the rows it
leaves out.

\vspace{0.5\baselineskip}
\begin{table}[H]
\caption{Prior art by layer. The last column states what Reproducible Science does at that layer; a row
marked \emph{below} is expanded in this appendix, and a section reference names where the body or
another appendix treats it. The claim-to-location row binds the unit Reproducible Science binds, a printed
claim to a location inside a cited artifact. The other rows describe a whole artifact, a run, an
event history or a semantic relation.}
\label{tab:related-layers}
\footnotesize
\begin{tabularx}{\linewidth}{@{}>{\raggedright\arraybackslash}p{0.21\linewidth}>{\raggedright\arraybackslash}p{0.31\linewidth}L@{}}
\toprule
Layer & Prior art & Relation to Reproducible Science \\
\midrule
entity, activity and agent provenance & PROV-O \citep{lebo2013provo}, OPM \citep{moreau2011opm}, PAV \citep{ciccarese2013pav} & defines no provenance ontology; exports an RO-Crate (below) \\
workflow structure & ProvONE \citep{cuevasvicenttin2016provone,cao2016provone} & ordering as a constraint in the plan; below \\
event ordering & PROV constraints \citep{cheney2013provconstraints} & consistency of a record with some ordering of its events; below \\
assertion-level publication & nanopublications \citep{groth2013nanopubguidelines,kuhn2016nanopub} & an assertion published as its own object; below \\
signed attestation & in-toto \citep{torresarias2019intoto}, SLSA \citep{openssf2023slsa}, Sigstore \citep{newman2022sigstore} & no report is signed, and a frozen plan is dated by OpenTimestamps (Appendix~\ref{sec:timestamps}); below \\
packaging at rest & RO-Crate \citep{soilandreyes2022rocrate}, Workflow Run RO-Crate \citep{leone2024wrroc} & export target: \code{repro export ro-crate}; below \\
input tracking & DVC \citep{iterative2020dvc}, AiiDA \citep{huber2020aiida}, tximeta \citep{love2020tximeta} & inputs pinned by digest, plus the address of each printed value; below \\
experiment tracking & Weights \& Biases \citep{wandb2026docs}, MLflow \citep{mlflow2026tracking} & a manuscript sentence bound to a file address, not a metric to a run; below \\
execution capture & ReproZip \citep{chirigati2016reprozip}, Binder \citep{jupyter2018binder} & which stored value backs which printed sentence; below \\
artifact-level verification & Code Ocean \citep{codeocean2026verification}, Whole Tale \citep{brinckman2019wholetale} & one outcome per assertion, not one per capsule or tale; Section~\ref{sec:related} \\
claim-to-evidence entailment & FEVER \citep{thorne2018fever}, SciFact \citep{wadden2020scifact}, CliniFact \citep{clinifact2025} & outcomes name a byte or decimal comparison, not a semantic label; Section~\ref{sec:related} \\
claim-to-work markup & ClaimReview \citep{schemaorg2026claimreview,google2026factcheck} & a claim addressed inside an artifact, not a whole work; below \\
claim-to-location binding & Paperclip \citep{gxl2026paperclip}, SciTeX Clew \citep{watanabe2026clew} & the same unit, compared deterministically against pinned bytes; Section~\ref{sec:related} \\
re-execution & DVC \citep{iterative2020dvc}, Snakemake \citep{koster2012snakemake}, CODECHECK \citep{nust2021codecheck}, ScientistOne \citep{meng2026scientistone}, ReScience C \citep{rougier2017rescience} & Appendix~\ref{sec:reexecution-precedents} \\
record linkage & Fellegi--Sunter \citep{fellegi1969linkage} & a third outcome for a decision not made; below \\
trace link recovery & information-retrieval traceability \citep{borg2013traceability} & links declared when the sentence is written; below \\
write-time binding & Sweave \citep{leisch2002sweave}, showyourwork \citep{luger2021mapping} & every printed value checked, generated or cited; below \\
repository auditing & adduce \citep{chudasama2026adduce} & adduce runs \reproverify{} as one of its rules; Section~\ref{sec:related} \\
agent trace graphs & LEDGER \citep{kim2026ledger} & checks declared values; reconstructs no trace; below \\
epistemic dependence & \citet{hardwig1985dependence}, \citet{wagenknecht2014opaque}, \citet{koskinen2023noepistemology}, \citet{ortmann2025oracles}, \citet{messeri2024illusions} & one link made recomputable; below \\
\bottomrule
\end{tabularx}
\end{table}

\paragraph{Provenance vocabularies.} PROV-O expresses the PROV Data Model in OWL2 as classes,
properties and restrictions for representing and interchanging provenance information, and three
classes form its basis: an entity, an activity that occurs over a period of time and acts upon
entities, and an agent bearing responsibility for an activity \citep{lebo2013provo}. PAV extends
PROV-O for the provenance, authoring and versioning of web resources, and distinguishes the
contributors, authors and curators of content from the creators of its representations
\citep{ciccarese2013pav}. Both vocabularies describe what happened to an artifact and who was
responsible for it. In Reproducible Science, a manifest entry records where in an artifact a printed value
sits, and the report records what a comparison at that address found.

\paragraph{Ordering in provenance standards.} ProvONE separates prospective provenance, the
workflow specification that produces a data artifact, from retrospective provenance, the steps
that were executed \citep{cuevasvicenttin2016provone}. On the specification side a
\code{Controller} states that one \code{Program}'s execution is controlled by another, so that in a
synchronous dataflow model a program starts only once a preceding program terminates; on the trace
side an \code{Execution} is associated with the \code{Program} it ran, which stands as its PROV plan
\citep{cuevasvicenttin2016provone}. DataONE adopted the model to collect, store, index and browse
the provenance of the data packages held on its member nodes \citep{cao2016provone}. OPM decorates
a causality graph with time and does not derive causality from it, because time precedence does
not imply causality: causal dependencies are asserted with edges, observed times serve to validate
them, and all clocks are assumed to be synchronized \citep{moreau2011opm}. PROV's ordering
constraints require a record to be consistent with a sensible ordering of its events, stated
through precedence relations, and for that check the times attached to events are irrelevant
\citep{cheney2013provconstraints}. Table~\ref{tab:ordering} sets the three standards beside the
check Reproducible Science makes.

\vspace{0.5\baselineskip}
\begin{table}[H]
\centering
\small
\caption{How three provenance standards represent ordering, beside the check Reproducible Science makes.
Each standard orders events within a single record. A confirmatory claim needs precedence
between two records, a registration and a run, with no causal edge asserted between them.}
\label{tab:ordering}
\begin{tabular}{@{}>{\raggedright\arraybackslash}p{0.17\textwidth}>{\raggedright\arraybackslash}p{0.33\textwidth}>{\raggedright\arraybackslash}p{0.42\textwidth}@{}}
\toprule
\textbf{Standard} & \textbf{How ordering appears} & \textbf{What it lets a reader conclude} \\
\midrule
OPM & observed times on a causality graph & that observed times agree with causal edges asserted separately \\
PROV constraints & precedence relations between events & that a record is consistent with some ordering of its events \\
ProvONE & control links in the workflow specification & which step must follow which in the plan \\
Reproducible Science & a run record read beside a pinned plan & that, as recorded, every run producing a confirmatory claim's evidence started after the plan's registration \\
\bottomrule
\end{tabular}
\end{table}

The check Reproducible Science makes rests on precedence alone. It does not assert that a registration caused a
run, and it reports a confirmatory claim \code{ordered}, \code{violated}, or \code{unchecked} with
a reason, such as a missing run record or a plan changed since it was pinned. A plan rewritten
after the results breaks its pin, and the claim's ordering reverts to \code{unchecked}, unless the
freeze record is rewritten with it. The registration time is self-recorded, but an outside timestamp bounds when the plan existed
(Section~\ref{sec:limits}, Appendix~\ref{sec:timestamps}).

\paragraph{Signed attestation.} In in-toto, whoever performs a step writes signed link metadata
recording that the step took place, and verification checks that each step defined has enough
signed links \citep{torresarias2019intoto}. Verifying an
artifact under SLSA means checking that a provenance statement's subject matches the artifact's
digest and that the provenance matches expected values \citep{openssf2023slsa}. Sigstore ties
signatures to existing developer identities through OIDC, signs with ephemeral keys, and records
signatures in artifact and identity logs \citep{newman2022sigstore}. Reproducible Science signs no report.
A freeze sends the plan's digest, and nothing else, to the OpenTimestamps calendar servers
(Appendix~\ref{sec:timestamps}), so the proof names no signer, which a plan frozen during
double-blind review requires.

\paragraph{Packaging and content-addressed inputs.} RO-Crate packages research artifacts with their
metadata as Schema.org annotations in JSON-LD, in a structured archive of the items that
contributed to a research outcome \citep{soilandreyes2022rocrate}, and Workflow Run RO-Crate is a
set of profiles that represent and package the provenance of a workflow run
\citep{leone2024wrroc}. tximeta identifies the reference transcriptome of a
quantification by a hashed checksum stored in the quantification output, and attaches annotation
metadata when the file is imported \citep{love2020tximeta}. AiiDA records every calculation or
workflow its engine runs in a provenance graph \citep{huber2020aiida}. Reproducible Science also pins its inputs
by digest, and adds the address of each value a manuscript prints from an output. \code{repro export ro-crate} writes our check records as an
RO-Crate: each claim as a Schema.org \code{Claim}, each checked assertion as a \code{CheckAction},
and the latest rerun of a declared command as a \code{CreateAction}, in which case the crate also
declares Process Run Crate. The fields for the address, the printed and stored values, and the
outcome are our own additions, defined in a profile we publish; RO-Crate has no terms for them.

\paragraph{Experiment tracking.} A Weights \& Biases run is a record of a single computation, to
which metrics, artifacts and other information are logged, and its artifacts track and version
data as the inputs and outputs of runs \citep{wandb2026docs}. MLflow Tracking logs parameters, code
versions, metrics and output files, and each run records its metadata and artifacts
\citep{mlflow2026tracking}. Both organize results by run. Reproducible Science ties a sentence of the
paper to a pinned file, which can also point to a Weights \& Biases run, an MLflow run, or another
external record.

\paragraph{Execution capture.} ReproZip traces one execution of an experiment and creates a
self-contained package of all its required dependencies, which unpacks into a directory, a chroot,
a virtual machine or a container \citep{chirigati2016reprozip}. Binder creates sharable,
interactive environments in the cloud from existing repositories, capturing a repository and its
environment in a container \citep{jupyter2018binder}. Capturing the environment shows that the
code can run again, but not which stored number supports which printed sentence. Reproducible Science records
that link as an assertion.

\paragraph{Replication by reimplementation.} ReScience C is a peer-reviewed journal for computational
research that encourages the explicit replication of published research through new, open-source
implementations \citep{rougier2017rescience}. A replication there writes the analysis again.
\code{repro reproduce} reruns the authors' own declared commands and compares each printed number
with the value at its declared address in the output (Appendix~\ref{sec:reexecution-precedents}).

\paragraph{ClaimReview.} Schema.org defines ClaimReview as a fact-checking review of claims made or
reported in a creative work, and reports it in use on 1K to 10K domains
\citep{schemaorg2026claimreview}. Google's fact-check structured data requires three of its
properties: \code{claimReviewed}, a short summary of the claim; \code{reviewRating}, the
assessment; and \code{url}, the page hosting the fact check \citep{google2026factcheck}. The claim
therefore travels as text. A \code{Claim} records where it occurred through \code{appearance} and
\code{firstAppearance}, and both take a \code{CreativeWork} \citep{schemaorg2026claimreview}, so the
finest address the vocabulary offers is a whole work. Fact-checking projects use rating schemes that
differ, each documented by its publisher, and search results show only the textual rating
\citep{google2026factcheck}. Whether a review separates a claim found false from a claim a reviewer
could not settle is therefore decided by each publisher, in its own words. Agreement between the
markup and the page rests with the publisher: the guidelines state that ``there must not be any
mismatch between the structured data and page content'' \citep{google2026factcheck}, and the
guidelines give no way to check it. \code{repro export ro-crate} writes each claim as a Schema.org
\code{Claim} and records beside it the address inside the artifact and the outcome of the
comparison made there.

\paragraph{Nanopublications.} A nanopublication consists of an assertion, the provenance of the
assertion, and publication information about the nanopublication itself, each an RDF graph, linked
to the nanopublication through a head graph \citep{groth2013nanopubguidelines}. Nanopublications are stored in a decentralized server network
and identified by trusty URIs, each carrying a hash computed on the RDF content it names, which
makes a nanopublication and its whole reference tree verifiable and immutable, and servers only add
entries and never update them \citep{kuhn2016nanopub}. A trusty URI plays the role a pin plays in
Reproducible Science (Section~\ref{sec:construction}). The units differ: a nanopublication publishes an
assertion as a new citable object carrying its own provenance, and our evidence assertion
adds nothing to the artifact it names, recording a location inside an artifact the manuscript
already cites.

\paragraph{The evidence assertion as a contract.} \citet{meyer1992contract} states a routine's
obligations as a precondition that any call must satisfy, a postcondition the routine ensures in
return, and a class invariant that constrains all routines of a class. Our evidence
assertion carries all three. The precondition is that the manifest loads, so that every assertion's locator
has the form its kind requires; a manifest that does not load reports nothing. The postcondition is
that every assertion evaluated is reported with its three stages and exactly one of the five
flattened labels of Table~\ref{tab:stages}. The invariant holds across the four kinds of
Table~\ref{tab:contract}: a locator resolves to exactly one scalar or to none, two or more matches
are reported as ambiguous rather than resolved to the first, and no backend searches a file for the
printed number.

\paragraph{Record linkage and traceability.} \citet{fellegi1969linkage} decide each record pair as
a link, a non-link or a possible link, and a possible link requires expensive manual linkage.
Comparing every pair of records would outstrip the capacity of even the largest and fastest
computers, so they restrict the comparisons examined to a subspace of the pairs by blocking
\citep{fellegi1969linkage}. Information-retrieval traceability recovers links between software
artifacts in a similar way: candidate links are ranked by the similarity between artifacts, and an
engineer examines each and confirms it or not \citep{borg2013traceability}. Recovering a link
after the fact produces some incorrect links \citep{borg2013traceability}. In Reproducible Science, an evidence
assertion is written with the sentence it binds, so no candidate set is ranked and no pair is searched for, and
what remains is whether the value at the declared address matches.

\paragraph{Write-time binding.} Sweave and showyourwork write each value into the document from the
code that computes it, so the two match after every build \citep{leisch2002sweave,luger2021mapping}. That
holds only while the code still runs, and it often does not: \citet{piccolo2016tools} trace
irreproducible findings to how software is packaged, installed and executed, and to how analysis
steps are documented. Reproducible Science checks each printed value against the stored result file, so the check
works whether or not the code still runs.

\paragraph{Spreadsheets as a results format.} Results are often published as Excel files: a
screen of open-access genomics papers in PubMed Central from 2014 to 2020 found 32{,}841 with
supplementary Excel files, and gene name errors appeared in 30.9\% of the 11{,}117 that held gene
lists, from gene names converted to dates, to floating-point numbers and to five-digit internal
dates \citep{abeysooriya2021genenames}. The table assertion of Reproducible Science reads such files directly,
finding a cell by sheet, column and a selector that matches one row (Table~\ref{tab:contract}).

\paragraph{Agent trace graphs.} LEDGER is a tracing and review system that builds layered trace
graphs over observed agent sessions \citep{kim2026ledger}. Its authors observe that visibility
alone leaves reviewers to reconstruct which actions, artifacts and validation steps matter for a
conclusion, and its traces expose workflow decisions, artifact lineage, repair steps, validation
coverage and claim-support paths for evidence-centered audit \citep{kim2026ledger}. \reproverify{}
checks each declared value without a reviewer, but does not guarantee the right value was
declared.

\paragraph{Epistemic dependence.} \citet{hardwig1985dependence} holds that a rational believer can
never avoid some dependence on experts, and observes of a paper with 99 authors that many of them
``will not even know how a given number in the article was arrived at.'' \citet{wagenknecht2014opaque}
separates two forms of that dependence within a research team: it is opaque when the dependent
scientist lacks the expertise to carry out and assess a colleague's work, and translucent when she
has it. Neither form is transparent, because in experimental science evidential contributions ``can
be entirely transparent only to those who have created or eye-witnessed them''
\citep{wagenknecht2014opaque}. For research that uses AI, \citet{koskinen2023noepistemology}
argues that no one can take informed responsibility for how an essentially opaque system reaches
its output, so that its opacity ``breaks the networks of trust assumed in social epistemology.''
\citet{ortmann2025oracles} replies that relying on such a system requires reliability rather than
trust, and cites language models paired with proof checkers, where the reliability of the final
output does not depend on how often the model errs. \citet{messeri2024illusions} argue that
scientists who trust AI tools to overcome their own cognitive limits become susceptible to
illusions of understanding, and call for exposure to science and technology studies, social
epistemology and philosophy of science. Reproducible Science checks reporting the way Ortmann's
proof checker checks a proof: any reader can recompute the link between a printed value or quotation
and its pinned artifact, while the extractor adds a dependence of its own (Appendix~\ref{sec:suite}).
